% Arbitrary-field virtual-annulus/tropical-ray comparison candidate.
% Date: 27 August 2026.
% This packet is not integrated into the book and makes no novelty claim.

\begin{proposition}
\label{prop:p21p3-gf-virtual-annulus-bridge}
Let \(K\) be a complete nontrivially valued non-Archimedean field, let
\(Y\) be a good strict rig-smooth \(K\)-analytic curve, and let \(x\) be an
interior point of type \(2\).  Write \(C_x\) for the normal projective
integral residual curve at \(x\).  If \(b\) is a branch at \(x\), write
\(p_b\in C_x\) for its corresponding closed point.

For \(g\in\mathcal O_{Y,x}^{\times}\), let
\(\partial_b(g)\in\mathbf Z\) be its exponent in the rank-one quotient of
the refined branch value group.  More explicitly, choose
\(s_b\in\operatorname{Rat}(C_x)^{\times}\) with
\(\operatorname{ord}_{p_b}(s_b)=1\), lift \(s_b\) to a norm-one analytic
unit \(f_b\in\mathcal O_{Y,x}^{\times}\), and require
\begin{equation}
\label{eq:p21p3-gf-branch-order-definition}
 \overline{|g(x)|_b}
 =\overline{|f_b(x)|_b}^{\,\partial_b(g)}
 \quad\text{in}\quad
 |\mathcal H(x)^{\times}|_b/|\mathcal H(x)^{\times}|.
\end{equation}
Then \(\partial_b(g)\) is independent of \(s_b\) and its lift subject to
\(\operatorname{ord}_{p_b}(s_b)=1\), and
\begin{equation}
\label{eq:p21p3-gf-branch-order-residue}
 \partial_b(g)=\operatorname{ord}_{p_b}(\widetilde g)
\end{equation}
whenever \(g\) has norm one at \(x\) and \(\widetilde g\) denotes its
degree-one residue.

The following assertions hold.
\begin{enumerate}
\item\label{itm:p21p3-gf-bridge:ray}
Let \(g_1,\ldots,g_m\in\mathcal O_{Y,x}^{\times}\), and assume that one
of them, say \(g_m\), satisfies \(\partial_b(g_m)=1\).  Put
\[
 G_{\mathrm{trop}}
   =(-\log|g_1|,\ldots,-\log|g_m|),
 \qquad
 v_b=(\partial_b(g_1),\ldots,\partial_b(g_m)).
\]
There are a cofinal virtual-annular section \(Z_b\) of \(b\) and a number
\(\varepsilon_b>0\) such that every \(|g_j|\) factors through the canonical
retraction \(\rho_b:Z_b\to S^{\mathrm{an}}(Z_b)\), the restriction of
\(G_{\mathrm{trop}}\) to \(S^{\mathrm{an}}(Z_b)\) is a homeomorphism, and
\begin{equation}
\label{eq:p21p3-gf-exact-branch-ray}
 G_{\mathrm{trop}}(Z_b)
 =G_{\mathrm{trop}}(x)+(0,\varepsilon_b)v_b.
\end{equation}
In particular, \(v_b\in\mathbf Z^m\) is the primitive outgoing tangent of
this Euclidean ray.

\item\label{itm:p21p3-gf-bridge:separation}
Let \(b_1,\ldots,b_s\) be distinct branches at \(x\), and let \(F_0\) be
any moment defined near \(x\).  There is a finite family of analytic units
\(f_1,\ldots,f_N\in\mathcal O_{Y,x}^{\times}\) such that, for the
concatenated moment \(F=(F_0,f_1,\ldots,f_N)\), the ray germs associated
with \(b_1,\ldots,b_s\) are pairwise distinct primitive rays and no branch
outside \(\{b_1,\ldots,b_s\}\) has the same outgoing ray germ as any
\(b_i\).

\item\label{itm:p21p3-gf-bridge:localization}
Let \(T\) be a positive \(\mathrm d'\)- and \(\mathrm d''\)-closed CLD
current of bidimension \((1,1)\), equivalently of bidegree \((0,0)\), near
\(x\).  In a source-separated
localization of a compactly supported test summand as in the proof of
\cite[Theorem~16.8.5]{CLD25}, replace its moment by the moment \(F\) from
\cref{itm:p21p3-gf-bridge:separation}.  For
\[
 \Omega=\mathbf R^{M}\setminus F_{\mathrm{trop}}(\partial V),
 \qquad
 U=F_{\mathrm{trop}}^{-1}(\Omega),
\]
the one-cell of the localized one-dimensional tropical image issuing from
\(F_{\mathrm{trop}}(x)\) in the \(b_i\)-direction is exactly the image of a
cofinal virtual-annular section of \(b_i\), and its primitive tangent is the
branch-order vector from
\cref{itm:p21p3-gf-bridge:ray}.  Thus, after a residue function
\(g\in\operatorname{Rat}(C_x)^\times\) and a norm-one analytic lift have
also been appended to the moment, the lift coordinate of the primitive
tangent is \(\operatorname{ord}_{p_{b_i}}(g)\); if the edge flux is
\(a_i\), its term in that scalar component of current balancing is
\(a_i\operatorname{ord}_{p_{b_i}}(g)\).  Here source-separated means that the
moment already contains the plateau tropicalizations used to distinguish
the finitely many source pieces meeting the test support.  This comparison
takes place entirely over \(K\); no pullback or descent of the current is
involved.
\end{enumerate}
\end{proposition}

\begin{remark}
\label{rem:p21p3-gf-virtual-annulus-index}
After a finite separable extension splitting \(Z_b\), a source annulus
coordinate can map to the primitive target coordinate with a positive
integer expansion index \(e_b\).  The proof shows that all coordinate
exponents are multiplied by the same \(e_b\), while the saturated target
tangent remains \(v_b\).  This is precisely why ordinary-current flux and
canonical calibrated flux acquire the same factor under tropical
pushforward.
\end{remark}
